% Apartado 1: reto y enfoque.
\section{Reto y enfoque}\label{sec:reto}

Unia es un asistente bancario conversacional, por voz y por texto, que trabaja sobre una base de datos ficticia
propia. Responde preguntas libres sobre los movimientos, dibuja gráficos y prepara Bizum, que es lo que pedía el reto.
También bloquea y desbloquea tarjetas. Esa función no estaba en el enunciado. La añadí al final, con el resto ya
terminado, como prueba de lo sencillo que es meter una funcionalidad nueva con esta arquitectura, y en el apartado~\ref{sec:escalabilidad}
cuento qué costó exactamente.

Al diseñarlo pensé en cuatro tipos de cliente. El cliente digital con prisa quiere preguntar ``¿cuánto gasté en
supermercados este mes?'' sin navegar por menús. Una persona mayor agradece hablar con normalidad, oír respuestas cortas
en voz alta y poder bloquear la tarjeta diciendo ``he perdido la tarjeta''. Una persona ciega o con baja visión
puede usarlo con voz de entrada y de salida, oír el código OTP cifra a cifra y aprobar con una huella (simulada).
La interfaz también está preparada para los lectores de pantalla (el cursor salta solo a la casilla del código y los SMS
se anuncian), aunque no la he podido probar con uno real.

Antes de entrar en las piezas conviene explicar la regla que lo ordena todo. Un modelo de lenguaje es muy bueno
entendiendo lo que pide una persona, pero también se equivoca, inventa cifras y se le puede engañar con un mensaje
preparado. En un banco eso no es admisible, así que repartí los papeles desde el principio: el modelo propone y el
servidor decide. El modelo solo interpreta la petición y elige qué herramienta llamar y con qué argumentos (por ejemplo,
``transferencia a Dani de 20 euros''). Lo que pasa a partir de ahí es código normal, sin modelo por medio: el servidor
comprueba el saldo y los topes, usa siempre el cliente que ha iniciado sesión y no el que diga el texto, y para mover
dinero exige un código de confirmación que el modelo nunca ve. Por eso, aunque alguien consiga que el modelo diga lo
que quiera, lo más que obtiene es una operación pendiente que nadie ha confirmado. Casi todo el diseño que sigue sale
de esta regla.

El asistente tiene ocho herramientas, definidas en \ruta{backend/aplicacion/herramientas/} y registradas en
\ruta{backend/composicion.py}. Cada una declara qué puede hacer con la cuenta del cliente: solo leer, escribir (mover
dinero o cambiar algo, y entonces exige el código de confirmación) o proteger (cambiar algo en la dirección segura, como
bloquear una tarjeta, sin código).

{\small
\begin{tabularx}{\textwidth}{@{}P{4.1cm}LP{3.1cm}@{}}
\toprule
\textbf{Herramienta} & \textbf{Qué hace} & \textbf{Riesgo declarado} \\
\midrule
\ruta{consultar_saldo} & Saldo, IBAN y titular & lectura \\
\ruta{listar_contactos_bizum} & Agenda Bizum del cliente & lectura \\
\ruta{consultar_movimientos} & Pregunta libre → SQL de solo lectura (NL→SQL) & lectura \\
\ruta{generar_grafico} & Gráfico calculado por el servidor o configuración Chart.js libre & lectura \\
\ruta{transferencia_bizum} & \textbf{Prepara} un Bizum: operación pendiente + código al móvil & escritura (exige OTP) \\
\midrule
\multicolumn{3}{@{}l@{}}{\textit{Añadidas al final como demostración de extensibilidad (apartado~\ref{sec:prueba-tarjeta}):}} \\
\ruta{consultar_tarjetas} & Tarjetas del cliente y su estado & lectura \\
\ruta{bloquear_tarjeta} & Bloquea al instante, sin código, auditado & protección \\
\ruta{desbloquear_tarjeta} & \textbf{Prepara} el desbloqueo: operación pendiente + código & escritura (exige OTP) \\
\bottomrule
\end{tabularx}}

La base de datos tiene cuatro clientes con vidas distintas (\ruta{backend/adaptadores/salida/sqlite/perfiles.py}). Alex
es un trabajador joven, Ana es estudiante, Marta tiene una familia con hijos y José Luis está jubilado. Cada uno tiene
ingresos, recibos, hábitos de gasto, contactos Bizum, una o dos tarjetas y dos años de movimientos.

La tabla siguiente indica en qué apartado trato cada requisito del reto.

{\small
\begin{tabularx}{\textwidth}{@{}LP{5.2cm}@{}}
\toprule
\textbf{Criterio} & \textbf{Apartado} \\
\midrule
Conversación cómoda y fluida, por voz y por texto, con poca latencia & \ref{sec:agente} (streaming, guardas) y \ref{sec:voz} (voz) \\
Operaciones: saldo y Bizum, con seguridad simulada & \ref{sec:seguridad} (Bizum con OTP, bloqueo de tarjeta) \\
Consultas en lenguaje natural sobre los datos (NL→SQL) & \ref{sec:nlsql} \\
Gráficos dinámicos: la IA decide cuál, sin plantillas & \ref{sec:graficos} \\
Refuerzo visual en pantalla & \ref{sec:comunicacion} (eventos SSE \cod{grafico}, \cod{tarjeta}, \cod{operacion}) \\
``Si mañana son 40 funciones, ¿hay que rehacer el agente?'' & \ref{sec:escalabilidad} (con el bloqueo de tarjeta como prueba real) \\
BD ficticia propia con varios clientes & \ref{sec:reto} y \ref{sec:seguridad} (aislamiento) \\
Accesibilidad & \ref{sec:voz} \\
\bottomrule
\end{tabularx}\par}
